\documentclass{article}

\usepackage{bm}
\usepackage{ctex}
\usepackage{tikz}
\usepackage{float}
\usepackage{xcolor}
\usepackage{setspace}
\usepackage{datetime}
\usepackage{geometry}
\usepackage{graphicx}
\usepackage{enumitem}
\usepackage{listings}
\usepackage{tcolorbox}
\usepackage{nicematrix}
\usepackage{amsmath, amssymb, amsfonts, mathrsfs}
\usepackage[fontsize = 12pt]{fontsize}

\renewcommand{\thesubsection}{(\alph{subsection})}
\newcommand{\diff}[2]{\dfrac{\mathrm{d}#1}{\mathrm{d}#2}}
\newcommand{\diffn}[3]{\dfrac{\mathrm{d}^{#3}#1}{\mathrm{d}#2^{#3}}}
\newcommand{\piff}[2]{\dfrac{\partial{#1}}{\partial{#2}}}
\newcommand{\eval}[2]{\left.#1\right|_{#2}}
\newcommand{\e}{\mathrm{e}}
\renewcommand{\d}{\mathrm{d}}
\newcommand{\barline}[1]{\overline{#1\rule{0pt}{1.9ex}}}

\geometry{
    a4paper,
    left = 2cm,
    right = 2cm,
    top = 2.4cm,
    bottom = 2.4cm
}

\setitemize[1]{
    itemsep = 7pt,
    partopsep = 0pt,
    parsep = \parskip,
    topsep = 5pt
}

\title{\vspace{-3em}应用统计：第三次作业}
\author{Illusionna \quad 2025.........004 \quad 人工智能学院}
\date{\currenttime,\ \today}



\begin{document}

\maketitle\vspace*{-1.5em}



\section{Q-2.2}
\textit{\textcolor{gray}{(1) 计算一组数据的均值、方差、标准差及均值的方差时，若给每个数据加上或减去一个正实数 $C$，与原数据组相比较，新数据组的均值、方差、标准差及均值的方差有无变化？为什么？}}

\textit{\textcolor{gray}{(2) 如果给每个原始数据乘以一个大于 1 的实数，结果又如何？有变化的请写出表达式。}}

为了方便表达，我们先设原数据组为 $x_1,\ x_2,\ \ldots,\ x_n$，则：

\begin{itemize}
    \item 均值
    \begin{itemize}
        \item 加减（有变化）：$\barline{x}_{\text{new}}=\dfrac{1}{n}\displaystyle\sum\limits_{i=1}^{n}\left(x_i\pm C\right)=\barline{x}\pm C$
        \item 数乘（有变化）：$\barline{x}_{\text{new}}=k\barline{x}$
    \end{itemize}
    \item 方差
    \begin{itemize}
        \item 加减（无变化）：$\sigma_{\text{new}}^2=\dfrac{1}{n}\displaystyle\sum\limits_{i=1}^{n}\left[\left(x_i\pm C\right)-\left(\barline{x}\pm C\right)\right]^2=\dfrac{1}{n}\displaystyle\sum\limits_{i=1}^{n}\left(x_i-\barline{x}\right)^2=\sigma^2$
        \item 数乘（有变化）：$\sigma_{\text{new}}^2=\dfrac{1}{n}\displaystyle\sum\limits_{i=1}^{n}\left(kx_i-k\barline{x}\right)^2=k^2\sigma^2$
    \end{itemize}
    \item 标准差
    \begin{itemize}
        \item 加减（无变化）：$\sigma_{\text{new}}=\sigma$
        \item 数乘（有变化）：$\sigma_{\text{new}}=\displaystyle\sqrt{\displaystyle k^2\sigma^2}=k\sigma$
    \end{itemize}
    \item 均值的方差
    \begin{itemize}
        \item 加减（无变化）：$D_{\text{new}}=\dfrac{\sigma^2}{n}$
        \item 数乘（有变化）：$D_{\text{new}}=\dfrac{k^2\sigma^2}{n}$
    \end{itemize}
\end{itemize}



\section{Q-2.3}
\textit{\textcolor{gray}{某厂制作一种产品，某天统计 10 名工人当天制作的数量，结果为 31,\ 23,\ 27,\ 30,\ 31,\ 32,\ 19,\ 26,\ 40,\ 31。构造该组数据的茎叶图，并求这组数据的平均值、中位数、众数、上四分位数、下四分位、极差、方差、标准差。}}

\[ \text{排序后数据}=[19,\ 23,\ 26,\ 27,\ 30,\ 31,\ 31,\ 31,\ 32,\ 40] \]

\begingroup
\linespread{1}\selectfont
\begin{lstlisting}[
    basicstyle = \ttfamily,
    columns = flexible,
    breaklines = true,
    escapeinside = {!@}{@!}
]
>>> stem | leaf
>>> -----------
>>>    1      9
>>>    2      3 6 7
>>>    3      0 1 1 1 2
>>>    4      0
\end{lstlisting}
\endgroup

\[ \barline{x}=\dfrac{1}{n}\sum_{i=1}^{n}x_i=\dfrac{290}{10}=29\qquad\text{median}=\dfrac{30+31}{2}=30.5 \]
\[ \text{mode}=31\qquad\text{range}=\max-\min=40-19=21 \]
\[ Q_1=26\qquad Q_3=31 \]
\[ S^2=\dfrac{1}{n-1}\sum_{i=1}^{n}\left(x_i-\barline{x}\right)^2=\dfrac{292}{10-1}\approx32.444\implies S=\sqrt{\dfrac{292}{9}}\approx5.696 \]



\section{Q-2.4}
\textit{\textcolor{gray}{设 $X\sim N(1,\ 9)$，从总体中抽取样本容量为 81 的样本，(1) 求样本均值 $\barline{X}$ 的标准差，(2) 求概率 $\Pr(0\leqslant\barline{X}\leqslant2)$.}}

根据总体分布 $X\sim N(1,\ 9)$，有：

\begin{itemize}
    \item 总体均值 $\mu=1$
    \item 总体方差 $\sigma^2=9$
    \item 总体标准差 $\sigma=3$
    \item 样本容量 $n=81$
\end{itemize}
\[ \sigma_{\overline{X}}=\dfrac{\sigma}{\sqrt{n}}=\dfrac{3}{\sqrt{81}}=\dfrac{1}{3} \]

\[ \barline{X}\sim N\left(\mu,\ \dfrac{\sigma^2}{n}\right)=N\left[1, \left(\dfrac{1}{3}\right)^2\right] \]

为了求概率，将 $\barline{X}$ 标准化，转化为标准正态分布 $Z\sim N(0,\ 1)$，标准化公式为：
\[ Z=\dfrac{\barline{X}-\mu}{\sigma_{\overline{X}}}=\dfrac{\barline{X}-1}{1/3} \]

当 $\barline{X}=0$ 时，$Z=-3$；当 $\barline{X}=2$ 时，$Z=3$，则根据 $3\sigma$ 准则有：
\[ \Pr(0\leqslant\barline{X}\leqslant2)=\Pr(-3\leqslant\barline{X}\leqslant3)=\Phi(3)-\Phi(-3)\approx99.73\% \]



\end{document}